% Wersja źródłowa: audyt II, jednostka 22
\setcounter{section}{21}
\section{Dystrybucje temperowane i transformacja Fouriera}
\label{sec:temperowane22}

\noindent\textbf{Fakty z wcześniejszych rozdziałów.}
Przyjmujemy, że przestrzeń Schwartza jest zupełną przestrzenią Frécheta,
$\mathcal D(\R^d)$ jest w niej gęsta, a transformacja Fouriera
jest jej ciągłym automorfizmem z odwrotnością i wzorem
$\mathcal F^2\varphi=\varphi(-\cdot)$.
Dystrybucja o zwartym nośniku $K$ działa również na funkcjach
$C^\infty$ poprzez $\langle T,f\rangle=\langle T,\chi f\rangle$,
gdzie $\chi=1$ w otoczeniu $K$, i spełnia oszacowanie
skończonego rzędu na ustalonym zwartym otoczeniu $K$.
Korzystamy z Twierdzeń~\ref{thm:schwartz-frechet17},
\ref{thm:density-schwartz17}, \ref{thm:fourier-schwartz17}
i Lematu~\ref{lem:compact-action21}.

W całym rozdziale $d\ge1$ jest liczbą całkowitą.
Rozważamy funkcje oraz dystrybucje zespolone na $\R^d$.
Przestrzeń Schwartza $\mathcal S=\mathcal S(\R^d;\C)$
wyposażamy w topologię Frécheta z rozdziału~17.
Dla $N,m\in\N_0$ przyjmujemy seminormy
$Q_{N,m}(\varphi)=\max_{|\alpha|\le m}\sup_{x\in\R^d}
(1+|x|)^N|\partial^\alpha\varphi(x)|$.
Przez $\mathcal S'$ rozumiemy przestrzeń \emph{ciągłych,
$\C$-liniowych} funkcjonałów na $\mathcal S$.
W parowaniu $\langle T,\varphi\rangle$ nie sprzęgamy zespolenie
funkcji testowej. Przypominamy, że
$\mathcal D(\R^d)=C_c^\infty(\R^d)$ jest gęsta w $\mathcal S$.

\subsection{Dystrybucje temperowane jako przestrzeń dualna}

\begin{definicja}[Dystrybucja temperowana]
\label{def:tempered22}
Dystrybucją temperowaną na $\R^d$ nazywamy element
$T\in\mathcal S'$, czyli ciągły liniowy funkcjonał
$T:\mathcal S\to\C$. Rozpatrujemy na $\mathcal S'$ topologię
słabą-* $\sigma(\mathcal S',\mathcal S)$, wyznaczoną przez
seminormy $T\mapsto|\langle T,\varphi\rangle|$
($\varphi\in\mathcal S$), oraz topologię silną
$\beta(\mathcal S',\mathcal S)$, wyznaczoną przez
$P_B(T)=\sup_{\varphi\in B}|\langle T,\varphi\rangle|$,
gdzie $B$ przebiega wszystkie zbiory ograniczone w $\mathcal S$.
\end{definicja}

\begin{twierdzenie}[Kryterium wzrostu seminormowego]
\label{thm:tempered-estimate22}
Niech $T:\mathcal S(\R^d;\C)\to\C$ będzie liniowy.
Następujące warunki są równoważne:
\begin{enumerate}
\item $T\in\mathcal S'(\R^d)$;
\item istnieją $C>0$ oraz $N,m\in\N_0$ takie, że
$|\langle T,\varphi\rangle|\le C Q_{N,m}(\varphi)$
dla wszystkich $\varphi\in\mathcal S$.
\end{enumerate}
W szczególności ograniczenie $T|_{\mathcal D}$ jest dystrybucją,
a odwzorowanie $\mathcal S'\to\mathcal D'$ dane przez
ograniczenie jest liniowe i różnowartościowe.
\end{twierdzenie}
\begin{proof}
Rodzina $Q_{j,j}$, $j\in\N_0$, jest rosnąca
i wyznacza topologię $\mathcal S$. Jeśli $T$ jest ciągły,
to dla pewnych $j$ i $\varepsilon>0$ warunek
$Q_{j,j}(\varphi)<\varepsilon$ implikuje
$|T(\varphi)|<1$. Liniowość i skalowanie dają
$|T(\varphi)|\le (2/\varepsilon)Q_{j,j}(\varphi)$;
gdy seminorma znika, wniosek wynika ze skalowania przez
dowolną liczbę. Odwrotna implikacja jest natychmiastowa.
Włożenie $j:\mathcal D\hookrightarrow\mathcal S$ jest ciągłe:
dla niepustego zwartego $K\subseteq\R^d$ i
$\varphi\in\mathcal D_K$ mamy
$Q_{N,m}(\varphi)\le (1+\max_{x\in K}|x|)^N p_m(\varphi)$.
Dla $K=\varnothing$ etap jest przestrzenią zerową.
Ciągłość $j$ wynika z własności uniwersalnej granicy
induktywnej (Twierdzenie~\ref{thm:LF-universal18}).
Stąd ograniczenie $T|_{\mathcal D}=T\circ j$ jest ciągłe.
Odwzorowanie $j^{\mathrm t}:T\mapsto T\circ j$ jest
liniowe i różnowartościowe: jeśli dwa ciągłe funkcjonały
zgadzają się na gęstej podprzestrzeni $\mathcal D$,
to ich różnica znika na całej $\mathcal S$.
\end{proof}

\begin{przyklad}[Dystrybucje regularne i miary]
\label{ex:tempered-regular22}
Jeżeli $f\in L^1_{\mathrm{loc}}(\R^d)$ i
dla pewnych $C>0$, $M\ge0$ zachodzi
$|f(x)|\le C(1+|x|)^M$ prawie wszędzie,
to $T_f(\varphi)=\int_{\R^d}f(x)\varphi(x)\,\mathrm d\lambda_d(x)$
określa element $\mathcal S'$. Istotnie, dla całkowitego $N>M+d$ mamy
\[
 \int_{\R^d}|f(x)\varphi(x)|\,\mathrm d\lambda_d(x)
 \le C Q_{N,0}(\varphi)
 \int_{\R^d}(1+|x|)^{M-N}\,\mathrm d\lambda_d(x).
\]
Ostatnia całka jest skończona: na kuli jednostkowej
funkcja jest ograniczona, a na pierścieniu
$2^k\le |x|<2^{k+1}$ całkę ogranicza stała razy
$2^{k(d+M-N)}$. Szereg tych oszacowań jest zbieżny.
Z Twierdzenia~\ref{thm:tempered-estimate22} wynika
temperowalność $T_f$.
Bardziej ogólnie, jeżeli $\mu$ jest zespoloną lokalnie skończoną
miarą borelowską i dla pewnego $N\ge0$ zachodzi
$\int_{\R^d}(1+|x|)^{-N}\,\mathrm d|\mu|(x)<\infty$,
to $\langle T_\mu,\varphi\rangle=
\int_{\R^d}\varphi(x)\,\mathrm d\mu(x)$
definiuje dystrybucję temperowaną. Istotnie, biorąc
liczbę całkowitą $n\ge N$, mamy
\[
 \left|\int_{\R^d}\varphi\,\mathrm d\mu\right|
 \le Q_{n,0}(\varphi)
 \int_{\R^d}(1+|x|)^{-N}\,\mathrm d|\mu|(x).
\]
Całka jest bezwzględnie zbieżna względem $|\mu|$,
a oszacowanie dowodzi ciągłości funkcjonału.
Przykładami są funkcja stała $1$, wielomiany i miary Diraca.
\end{przyklad}

\begin{przyklad}[Dystrybucja nietemperowana]
\label{ex:non-tempered22}
Dystrybucja regularna wyznaczona przez
$f(x)=e^{|x|^2}$ należy do $\mathcal D'(\R^d)$,
ale nie jest ograniczeniem żadnego elementu $\mathcal S'$.
Aby to wykazać, weźmy $\chi\in\mathcal D(\R^d)$,
$0\le\chi\le1$, równą $1$ na $\overline B(0,1)$,
i funkcje $\varphi_R(x)=\chi(x/R)e^{-|x|^2/2}$.
Dla $R\to\infty$ zachodzi
$\varphi_R\to e^{-|x|^2/2}$ w $\mathcal S$
na mocy dowodu Twierdzenia~\ref{thm:density-schwartz17}:
dla $g(x)=e^{-|x|^2/2}\in\mathcal S$ i $R\ge1$ zachodzi
$Q_{N,m}(g-\chi(\cdot/R)g)\le
C_m(1+R)^{-1}Q_{N+1,m}(g)$.
Przynależność $g$ do $\mathcal S$ wynika stąd, że
każda jego pochodna jest wielomianem pomnożonym przez $g$.
Jednak
$\langle T_f,\varphi_R\rangle
=\int_{\R^d}\chi(x/R)e^{|x|^2/2}\,\mathrm d\lambda_d(x)
\to+\infty$,
gdyż dla $R\ge1$ całka jest nie mniejsza niż
$\int_{B(0,R)}1\,\mathrm d\lambda_d=
R^d\lambda_d(B(0,1))$.
Gdyby istniało ciągłe przedłużenie na $\mathcal S$,
wartości te musiałyby zbiegać do skończonej liczby.
\end{przyklad}

\begin{stwierdzenie}[Podprzestrzenie klasycznych dystrybucji]
\label{prop:schwartz-embeddings22}
Istnieją naturalne, różnowartościowe i ciągłe odwzorowania
\[
 \mathcal D\hookrightarrow\mathcal S\hookrightarrow
 L^2(\R^d)\hookrightarrow\mathcal S'\hookrightarrow\mathcal D',
\]
przy czym obie przestrzenie dualne wyposażamy albo
w topologie silne, albo w topologie słabe-*,
gdzie funkcję $f\in L^2$ identyfikujemy z funkcjonałem
$\langle T_f,\varphi\rangle=\int_{\R^d}f(x)\varphi(x)
\,\mathrm d\lambda_d(x)$.
Ponadto każda dystrybucja o zwartym nośniku jest temperowana:
$\mathcal E'(\R^d)\subseteq\mathcal S'(\R^d)$.
\end{stwierdzenie}
\begin{proof}
Pierwsze włożenie wykazano w Twierdzeniu~
\ref{thm:tempered-estimate22}. Drugie jest ciągłe na mocy
Stwierdzenia~\ref{prop:schwartz-operators17} i różnowartościowe,
ponieważ ciągła funkcja równa zeru prawie wszędzie znika
wszędzie. Dla $f\in L^2$ nierówność Cauchy'ego--Schwarza daje
$|\langle T_f,\varphi\rangle|\le\|f\|_2\|\varphi\|_2$;
także $\|\varphi\|_2\le C_dQ_{d+1,0}(\varphi)$.
Dla ograniczonego $B\subseteq\mathcal S$ otrzymujemy
\[
 P_B(T_f)\le C_d\|f\|_2\sup_{\varphi\in B}Q_{d+1,0}(\varphi).
\]
Prawa strona jest skończona i dowodzi ciągłości silnej,
a dla singletonów także słabej-*.
Jeśli $T_f=0$, to $\int f\varphi\,\mathrm d\lambda_d=0$
dla wszystkich $\varphi\in\mathcal D$.
Z gęstości $\mathcal D$ w $L^2$ wybieramy $\varphi_n\to
\overline f$ w $L^2$; nierówność Cauchy'ego--Schwarza daje
$\int |f|^2\,\mathrm d\lambda_d=0$, więc $f=0$ w $L^2$.
Ostatnie odwzorowanie jest różnowartościowe z
Twierdzenia~\ref{thm:tempered-estimate22}. Dla
$B\subseteq\mathcal D$ ograniczonego mamy
$P_B(T|_{\mathcal D})=P_{j(B)}(T)$, gdzie $j(B)$ jest
ograniczony w $\mathcal S$ dzięki ciągłości $j$.
Daje to ciągłość silną, a oceny na pojedynczych
funkcjach testowych dają ciągłość słabą-*.
Jeżeli $T\in\mathcal E'$ ma zwarty nośnik,
to jego przedłużenie z Lematu~\ref{lem:compact-action21}
spełnia $|T(\varphi)|\le C Q_{0,m}(\varphi)$.
Jest więc ciągłe na $\mathcal S$ i zgadza się z $T$
na $\mathcal D$; identyfikujemy $T$ z tym jednoznacznym
przedłużeniem.
\end{proof}

\subsection{Operatory na przestrzeni dystrybucji temperowanych}

\begin{twierdzenie}[Różniczkowanie i mnożenie przez wielomian]
\label{thm:tempered-operators22}
Dla każdego $\alpha,\beta\in\N_0^d$ operatory
$T\mapsto\partial^\alpha T$ i $T\mapsto x^\beta T$
są liniowe i ciągłe na $\mathcal S'$ zarówno dla topologii
$\sigma(\mathcal S',\mathcal S)$, jak i
$\beta(\mathcal S',\mathcal S)$, gdzie
\[
\langle\partial^\alpha T,\varphi\rangle
=(-1)^{|\alpha|}\langle T,\partial^\alpha\varphi\rangle,
\qquad\langle x^\beta T,\varphi\rangle
=\langle T,x^\beta\varphi\rangle.
\]
To samo dotyczy mnożenia przez każdą funkcję
$a\in C^\infty(\R^d;\C)$, dla której dla każdego
$\gamma\in\N_0^d$ istnieją $C_\gamma>0$ oraz
$M_\gamma\in\N_0$ takie, że
$|\partial^\gamma a(x)|\le C_\gamma(1+|x|)^{M_\gamma}$
dla wszystkich $x\in\R^d$.
\end{twierdzenie}
\begin{proof}
Różniczkowanie i mnożenie przez $x^\beta$ są ciągłymi
operatorami liniowymi $\mathcal S\to\mathcal S$,
co wykazano w Stwierdzeniu~\ref{prop:schwartz-operators17}.
Dla ogólnego mnożnika $a$ ustalmy $N,m\in\N_0$,
$M_m=\max_{|\gamma|\le m}M_\gamma$ i
$C_m=\max_{|\beta|\le m}\sum_{\gamma\le\beta}
\binom\beta\gamma C_\gamma$.
Ze wzoru Leibniza, dla $|\beta|\le m$, dostajemy
\[
 (1+|x|)^N|\partial^\beta(a\varphi)(x)|
 \le C_m Q_{N+M_m,m}(\varphi).
\]
Po wzięciu supremum oraz maksimum otrzymujemy
$Q_{N,m}(a\varphi)\le C_m Q_{N+M_m,m}(\varphi)$.
Zatem odwzorowanie $\varphi\mapsto a\varphi$ jest
ciągłym operatorem na $\mathcal S$.
Dla dowolnego ciągłego operatora liniowego
$A:\mathcal S\to\mathcal S$ jego transpozycja
$A^{\mathrm t}T=T\circ A$ spełnia
$P_B(A^{\mathrm t}T)=P_{A(B)}(T)$.
Ciągłość $A$ zapewnia ograniczoność $A(B)$, skąd
ciągłość silna; dla singletonów otrzymujemy ciągłość
słabą-*. Stosujemy to do $A\varphi=(-1)^{|\alpha|}
\partial^\alpha\varphi$ i $A\varphi=a\varphi$,
zgodnie z definicjami pochodnej i mnożenia dystrybucji.
Ograniczenie tych działań do $\mathcal D$ zgadza się
z działaniami na zwykłych dystrybucjach.
\end{proof}

\begin{lemat}[Względna zwartość zbiorów ograniczonych w $\mathcal S$]
\label{lem:montel-schwartz22}
Każdy zbiór ograniczony w $\mathcal S(\R^d)$ ma zwarte
domknięcie w jej topologii Frécheta.
\end{lemat}
\begin{proof}
Niech $B\subseteq\mathcal S$ będzie ograniczony i
$(\varphi_n)\subseteq B$ dowolnym ciągiem.
Dla każdej liczby całkowitej $R\ge1$ i każdego $m\ge0$
rodziny pochodnych $\partial^\alpha\varphi_n$ dla
$|\alpha|\le m$ są jednostajnie ograniczone i
jednostajnie równociągłe na $\overline B(0,R)$:
równociągłość wynika z ograniczenia pochodnych
rzędu $m+1$ i całkowania wzdłuż odcinka w kuli:
dla $x,y$ w tej kuli oraz $|\alpha|\le m$ mamy
$|\partial^\alpha\varphi_n(x)-\partial^\alpha\varphi_n(y)|
\le\|x-y\|_1\sup_{\psi\in B}Q_{0,m+1}(\psi)$. Z twierdzenia Arzelà--Ascoliego
i argumentu diagonalnego uzyskujemy podciąg
$(\varphi_{n_k})$, którego wszystkie pochodne zbiegają
jednostajnie na wszystkich kulach domkniętych.
Nie musimy osobno identyfikować granic pochodnych:
poniższe oszacowanie da zbieżność w zupełnej
przestrzeni Schwartza.
Dla ustalonych $N,m$ i $R\ge1$ poza $B(0,R)$ zachodzi
\[
(1+|x|)^N|\partial^\alpha(\varphi_{n_k}-\varphi_{n_l})(x)|
\le R^{-1}\bigl(Q_{N+1,m}(\varphi_{n_k})
+Q_{N+1,m}(\varphi_{n_l})\bigr)\quad(|\alpha|\le m).
\]
Prawa strona jest jednostajnie mała przy dużym $R$;
wewnątrz kuli $B(0,R)$ zbieżność wszystkich pochodnych
jest jednostajna. Dlatego $(\varphi_{n_k})$ jest ciągiem
Cauchy'ego we wszystkich seminormach $Q_{N,m}$.
Zupełność $\mathcal S$ daje zbieżność do elementu
$\mathcal S$. Ponieważ przestrzeń jest metryzowalna,
ciągowa względna zwartość zbioru jest równoważna
zwartości jego domknięcia. Istotnie, dla ciągu $(g_n)$
w domknięciu $B$ wybieramy $f_n\in B$ w odległości
metrycznej mniejszej niż $1/n$ od $g_n$.
Zbieżny podciąg $(f_{n_k})$ daje zbieżny podciąg
$(g_{n_k})$ o tej samej granicy, należącej do domknięcia.
Zwartość wynika z kryterium ciągowego dla przestrzeni
metrycznych.
\end{proof}

\begin{twierdzenie}[Zbieżność słaba-* i silna ciągów]
\label{thm:seq-topologies-schwartz22}
Niech $(T_n)$ będzie ciągiem dystrybucji temperowanych,
a $T\in\mathcal S'$. Wtedy
\[
T_n\longrightarrow T\text{ w }\sigma(\mathcal S',\mathcal S)
\quad\Longleftrightarrow\quad
T_n\longrightarrow T\text{ w }\beta(\mathcal S',\mathcal S).
\]
Mimo to topologie $\sigma(\mathcal S',\mathcal S)$ i
$\beta(\mathcal S',\mathcal S)$ nie są równe.
\end{twierdzenie}
\begin{proof}
Implikacja ze zbieżności silnej jest natychmiastowa.
Załóżmy zbieżność słabą-*. Rodzina
$\{T_n-T:n\ge1\}$ jest punktowo ograniczona.
Ponieważ $\mathcal S$ jest przestrzenią Frécheta,
a więc beczkową, twierdzenie Banacha--Steinhausa
implikuje jej równociągłość.
Dla dowolnego ograniczonego $B\subseteq\mathcal S$
Lemat~\ref{lem:montel-schwartz22} daje zwartość
$K=\overline B$. Równociągłość funkcjonałów i
punktowa zbieżność do zera zapewniają zbieżność
jednostajną na $K$: dla danego $\varepsilon>0$
wybieramy otoczenie $V$ z
$\sup_n|\langle T_n-T,v\rangle|<\varepsilon/2$
dla $v\in V$ i skończone pokrycie $K$ przesunięciami
$\varphi_j+V$. Następnie dla dostatecznie dużego $n$
wszystkie $|\langle T_n-T,\varphi_j\rangle|<\varepsilon/2$.
Wówczas $P_B(T_n-T)<\varepsilon$.

Aby wykazać różność topologii, wybierzmy niezerową funkcję
$\psi\in\mathcal D$ o nośniku w $B(0,1)$ i zdefiniujmy
$\psi_n(x)=e^{-n^2}\psi(x-3ne_1)$ dla $n\ge1$.
Zbiór $B=\{\psi_n:n\ge1\}$ jest ograniczony w
$\mathcal S$: dla ustalonych $N,m$ mamy
$Q_{N,m}(\psi_n)\le e^{-n^2}(2+3n)^N Q_{0,m}(\psi)$,
co jest ograniczone niezależnie od $n$.
Jego liniowa otoczka jest nieskończenie wymiarowa,
ponieważ nośniki są parami rozłączne i każda $\psi_n$
jest niezerowa.
Otoczenie silne $\{T:P_B(T)<1/2\}$ nie zawiera
żadnego słabo-* otoczenia zera: dla dowolnie wielu
$\varphi_1,\ldots,\varphi_r\in\mathcal S$
można wybrać $\psi_n$ poza ich liniową otoczką.
Podprzestrzeń $M=\operatorname{span}\{\varphi_1,\ldots,
\varphi_r\}$ jest skończenie wymiarowa, a więc domknięta
w przestrzeni lokalnie wypukłej Hausdorffa $\mathcal S$.
Zatem $\mathcal S/M$ jest przestrzenią lokalnie wypukłą
Hausdorffa, a klasa $\psi_n+M$ jest niezerowa.
Twierdzenie Hahna--Banacha daje ciągły funkcjonał
na ilorazie niezerowy na tej klasie; po normalizacji
i złożeniu z odwzorowaniem ilorazowym otrzymujemy
$T\in\mathcal S'$ taki, że $T(\varphi_j)=0$ dla
wszystkich $j$ oraz $T(\psi_n)=1$.
Należy on do dowolnego podstawowego słabo-* otoczenia
wyznaczonego przez te $\varphi_j$, ale nie do
$\{T:P_B(T)<1/2\}$. Ponieważ topologia silna jest
co najmniej tak mocna jak słaba-*, dowodzi to ich różności.
\end{proof}

\subsection{Rozszerzenie transformacji Fouriera}

\begin{definicja}[Transformacja Fouriera dystrybucji temperowanej]
\label{def:fourier-tempered22}
Dla $T\in\mathcal S'(\R^d)$ definiujemy
$\widehat T=\mathcal F T\in\mathcal S'$ przez
\[
\langle\mathcal F T,\varphi\rangle
=\langle T,\mathcal F\varphi\rangle
\qquad(\varphi\in\mathcal S).
\]
Występuje tu transformacja Fouriera na $\mathcal S$
z normalizacją $(2\pi)^{-d/2}$ i jądrem
$e^{-ix\cdot\xi}$ z wcześniejszego rozdziału o przestrzeni Schwartza.
\end{definicja}

\begin{twierdzenie}[Transformacja Fouriera jest automorfizmem]
\label{thm:fourier-auto-tempered22}
Transformacja $\mathcal F:\mathcal S'\to\mathcal S'$
jest ciągłym izomorfizmem liniowo-topologicznym
zarówno dla topologii $\sigma(\mathcal S',\mathcal S)$,
jak i $\beta(\mathcal S',\mathcal S)$.
Jej odwrotność dana jest transpozycją
$\mathcal F^{-1}$ na $\mathcal S$. Zachodzą wzory
\[
\mathcal F(\partial_jT)=i\xi_j\,\mathcal F T,
\qquad
\mathcal F(x_jT)=i\partial_{\xi_j}(\mathcal F T),
\qquad\mathcal F^2T=T(-\cdot),
\]
gdzie $\langle T(-\cdot),\varphi\rangle=
\langle T,\varphi(-\cdot)\rangle$.
\end{twierdzenie}
\begin{proof}
Twierdzenie o transformacji Fouriera na przestrzeni Schwartza mówi, że
$\mathcal F:\mathcal S\to\mathcal S$ jest ciągłym
izomorfizmem liniowo-topologicznym, a jego odwrotność
także jest ciągła. Transponowanie obu operatorów
zapewnia wzajemnie odwrotne odwzorowania na $\mathcal S'$.
Ciągłość słaba-* wynika z działania na pojedynczej
funkcji testowej. Dla silnej mamy
$P_B(\mathcal FT)=P_{\mathcal F(B)}(T)$;
$\mathcal F(B)$ jest ograniczony, gdy $B$ jest ograniczony,
co daje ciągłość. Odwrotność rozpatrujemy analogicznie.
Dla $\varphi\in\mathcal S$ i współrzędnej $j$ otrzymujemy
\[
\begin{aligned}
\langle\mathcal F(\partial_jT),\varphi\rangle
&=-\langle T,\partial_j\mathcal F\varphi\rangle
=-\langle T,\mathcal F(-ix_j\varphi)\rangle
=\langle i\xi_j\mathcal FT,\varphi\rangle,\\
\langle\mathcal F(x_jT),\varphi\rangle
&=\langle T,x_j\mathcal F\varphi\rangle
=\langle T,\mathcal F(-i\partial_j\varphi)\rangle
=\langle i\partial_{\xi_j}\mathcal FT,\varphi\rangle.
\end{aligned}
\]
W tych rachunkach zmienna w funkcji testowej i
wektor w przestrzeni dualnej są identyfikowane
tylko notacyjnie; kluczowe są wzory
$\partial\mathcal F\varphi=\mathcal F(-ix\varphi)$
i $x\mathcal F\varphi=\mathcal F(-i\partial\varphi)$.
Ostatnia tożsamość wynika z
$\mathcal F^2\varphi(x)=\varphi(-x)$ i transpozycji.
\end{proof}

\begin{przyklad}[Funkcja stała, delta Diraca i pochodne]
\label{ex:fourier-delta22}
Dla $a\in\R^d$ mamy w $\mathcal S'$ wzory
\[
\mathcal F\delta_a=(2\pi)^{-d/2}e^{-ia\cdot\xi},
\quad\mathcal F T_1=(2\pi)^{d/2}\delta_0,
\quad\mathcal F(\partial^\alpha\delta_0)
=(2\pi)^{-d/2}(i\xi)^\alpha.
\]
Pierwszy wynika z
$\langle\mathcal F\delta_a,\varphi\rangle
=(2\pi)^{-d/2}\int e^{-ia\cdot\xi}\varphi(\xi)
\,\mathrm d\lambda_d(\xi)$.
Drugi można uzyskać, stosując $\mathcal F^2T=T(-\cdot)$
do pierwszego wzoru dla $a=0$. Trzeci wynika z
Twierdzenia~\ref{thm:fourier-auto-tempered22}.
W szczególności delta Diraca nie jest funkcją,
ale jej transformata Fouriera jest funkcją gładką.
\end{przyklad}

\begin{twierdzenie}[Twierdzenie Plancherela]
\label{thm:plancherel22}
Transformacja Fouriera z normalizacją $(2\pi)^{-d/2}$
przedłuża się jednoznacznie do operatora unitarnego
$\mathcal F:L^2(\R^d;\C)\to L^2(\R^d;\C)$,
czyli dla $f,g\in L^2$ zachodzi
$\int_{\R^d}\widehat f(\xi)\overline{\widehat g(\xi)}
\,\mathrm d\lambda_d(\xi)
=\int_{\R^d}f(x)\overline{g(x)}\,\mathrm d\lambda_d(x)$.
Dla $f\in L^2$ transformata rozumiana w $L^2$
zgadza się z transformatą dystrybucyjną $T_f$.
\end{twierdzenie}
\begin{proof}
Dla $f,g\in\mathcal S$ wzór odwrotny oraz Fubini
(dopuszczalny dzięki $f,\widehat g\in L^1$) dają
\[
\int\widehat f(\xi)\overline{\widehat g(\xi)}
\,\mathrm d\lambda_d(\xi)
=\int f(x)\overline{g(x)}\,\mathrm d\lambda_d(x),
\]
gdzie w obu całkach obszarem całkowania jest $\R^d$.
Zatem $\mathcal F$ jest izometrią $\mathcal S$ w $L^2$.
Przestrzeń $\mathcal S$ jest gęsta w $L^2$,
więc operator przedłuża się jednoznacznie do izometrii
na $L^2$. Jego obraz jest domknięty, a zawiera
$\mathcal F(\mathcal S)=\mathcal S$, zbiór gęsty;
stąd jest to operator unitarny. Tożsamość polaryzacyjna
przedłuża zachowanie iloczynu skalarnego na całe $L^2$.
Jeśli $f_n\in\mathcal S$ i $f_n\to f$ w $L^2$,
to $T_{f_n}\to T_f$ w $\mathcal S'$, ponieważ
$|\langle T_{f_n-f},\varphi\rangle|
\le\|f_n-f\|_2\|\varphi\|_2$.
Ciągłość transformaty w $\mathcal S'$ i jej zgodność
z klasyczną na $\mathcal S$ zapewniają ostatnią tezę.
\end{proof}

\subsection{Splot z funkcją szybko malejącą i transformata dystrybucji o zwartym nośniku}

\begin{twierdzenie}[Splot z elementem przestrzeni Schwartza]
\label{thm:schwartz-convolution22}
Niech $T\in\mathcal S'(\R^d)$ i
$\rho\in\mathcal S(\R^d)$. Funkcja
$u(x)=\langle T_y,\rho(x-y)\rangle$ należy do
$C^\infty(\R^d)$, a każda jej pochodna ma
co najwyżej wzrost wielomianowy. Dystrybucja
regularna $T_u$ jest temperowana i oznaczamy ją $T*\rho$.
Zachodzi
\[
\mathcal F(T*\rho)=(2\pi)^{d/2}(\mathcal FT)\,(\mathcal F\rho),
\qquad \partial^\alpha(T*\rho)=T*(\partial^\alpha\rho).
\]
\end{twierdzenie}
\begin{proof}
Dla $x$ w zbiorze zwartym odwzorowanie
$x\mapsto\rho(x-\cdot)$ jest gładkie jako funkcja
z wartościami w $\mathcal S$; jego pochodne to
$x\mapsto(\partial^\alpha\rho)(x-\cdot)$.
Zatem można różniczkować pod działaniem $T$.
Z Twierdzenia~\ref{thm:tempered-estimate22} mamy
$|T(\varphi)|\le C Q_{N,m}(\varphi)$.
Wstawiając $\varphi(y)=\partial^\alpha\rho(x-y)$
i korzystając z
$1+|y|\le(1+|x|)(1+|x-y|)$,
otrzymujemy
$|\partial^\alpha u(x)|\le C_\alpha(1+|x|)^N$.
Stąd $u$ jest temperowana jako funkcja.
Dla $\psi\in\mathcal S$ przez zamianę działania $T$
i całkowania funkcji $\psi(x)\rho(x-y)$
(uzasadnioną przez zbieżność w seminormach Schwartza)
otrzymujemy
$\langle T_u,\psi\rangle=
\langle T,\check\rho*\psi\rangle$,
gdzie $\check\rho(x)=\rho(-x)$.
To także dowodzi ciągłości konstrukcji w $\mathcal S'$.
Z twierdzenia splotowego w $\mathcal S$ wynika
$\check\rho*\mathcal F\psi
=(2\pi)^{d/2}\mathcal F((\mathcal F\rho)\psi)$.
Zatem
\[
\begin{aligned}
\langle\mathcal F(T*\rho),\psi\rangle
&=\langle T,\check\rho*\mathcal F\psi\rangle\\
&=(2\pi)^{d/2}\langle\mathcal FT,
 (\mathcal F\rho)\psi\rangle,
\end{aligned}
\]
co daje pierwszy wzór. Drugi wynika z różniczkowania
pod działaniem dystrybucji.
\end{proof}

\begin{twierdzenie}[Transformata dystrybucji o zwartym nośniku]
\label{thm:compact-fourier22}
Jeżeli $T\in\mathcal E'(\R^d)$ ma zwarty nośnik $K$,
to jego transformata Fouriera jest dystrybucją regularną
reprezentowaną przez gładką funkcję
\[
F_T(\xi)=(2\pi)^{-d/2}\langle T_x,e^{-ix\cdot\xi}\rangle.
\]
Dla pewnego $m\in\N_0$ i każdego wielowskaźnika $\beta$
istnieje $C_\beta>0$ takie, że
$|\partial_\xi^\beta F_T(\xi)|
\le C_\beta(1+|\xi|)^m$
dla wszystkich $\xi\in\R^d$.
W działaniu na funkcji wykładniczej używamy
przedłużenia dystrybucji o zwartym nośniku
z lematu o przedłużaniu dystrybucji o zwartym nośniku.
\end{twierdzenie}
\begin{proof}
Wybierzmy $\chi\in\mathcal D$ równą $1$ na
otoczeniu $K$. Dla każdej $\beta$ możemy różniczkować
$x\mapsto\chi(x)(-ix)^\beta e^{-ix\cdot\xi}$
po parametrze $\xi$, gdyż funkcje mają wspólny zwarty nośnik.
Dlatego
$\partial_\xi^\beta F_T(\xi)
=(2\pi)^{-d/2}\langle T,(-ix)^\beta e^{-ix\cdot\xi}\rangle$.
Oszacowanie z lematu o przedłużaniu dystrybucji o zwartym nośniku
przy ustalonym rzędzie $m$ i zwartym otoczeniu nośnika
$T$, razem ze wzorem Leibniza, daje
$|\partial_\xi^\beta F_T(\xi)|
\le C_\beta(1+|\xi|)^m$.
Jeśli $\varphi\in\mathcal S$, to
$\mathcal F\varphi\in\mathcal S$; zamiana działania
$T$ i całkowania szybko malejącego parametru jest
dozwolona z tego samego oszacowania. Otrzymujemy
\[
\langle\mathcal FT,\varphi\rangle
=\int_{\R^d}F_T(\xi)\varphi(\xi)
\,\mathrm d\lambda_d(\xi),
\]
co identyfikuje dystrybucję.
\end{proof}

\begin{uwaga}[Zakres rozdziału]
\label{rem:scope22}
Transformacja Fouriera jest tu automorfizmem przestrzeni
dystrybucji temperowanych i podstawowym narzędziem teorii
operatorów. Jej zastosowania do rozwiązywania konkretnych
równań różniczkowych pozostawiamy osobnemu skryptowi.
Dalsze rozdziały rozwiną teorię przestrzeni Montela,
przestrzeni nuklearnych i produktów tensorowych topologicznych.
\end{uwaga}
